\section{A tree of maps and separating functionals}
\label{sec:tree}

Fix the point $x$ supplied by Proposition~\ref{prop:anchor} and the
sequence $b_j=R_j(x)$ from Section~\ref{sec:anchor}.  We construct a
single infinite tree.  Each node carries a map in $\mathcal G$, and
each nonroot node also carries a functional.  The maps will allow a
subtree to approximate its prescribed label or to transmit a common
point $b_j$.  Each functional must be chosen before the weights in its own
sibling group are assigned.  A family of possible root outputs parameterized
by a compact finite-dimensional set makes this order of choices possible.

\subsection{Evaluation on an ordered tree}

The root is denoted by $r$, and $L_k$ denotes the finite, nonempty
level of depth $k$, with $L_0=\{r\}$.  Every node has a finite,
nonempty ordered list of children.  The order on $L_{k+1}$ is obtained
by placing these lists consecutively in the order of their parents
in $L_k$.  Write $\widehat v$ for the parent of $v$, and $p\prec v$
when $p$ is a proper ancestor of $v$.  In particular, $\prec$ does
not denote the order within a level.

Each node $v$ is assigned a map $G_v\in\mathcal G$ and the point
$y_v=G_v(x)$.  At the root, $G_r=\operatorname{Id}_K$ and $y_r=x$.
Each nonroot node has a weight $\alpha_v>0$, with
\[
  \sum_{\widehat v=u}\alpha_v=1
  \qquad\text{for every node }u.
\]
Given outputs at the children of $u$, the output at $u$ is
\begin{equation}
  H_u=G_u\left(\sum_{\widehat v=u}\alpha_v H_v\right).
  \label{eq:tree-rule}
\end{equation}
Every output belongs to $K$.  If outputs $(H_v)_{v\in L_k}$ are
prescribed at a cutoff level, let $\Phi_k((H_v)_{v\in L_k})$ be the
root output obtained by applying this rule strictly above that level.
Thus $\Phi_0(H_r)=H_r$; a map at the cutoff itself is not applied
again.

Define the path weights by
\[
  W_r=1,\qquad W_v=W_{\widehat v}\alpha_v\quad(v\ne r).
\]
An induction gives $\sum_{v\in L_k}W_v=1$.  More generally, the
sum of the weights at any later level below a given node is the
weight of that node.  Nonexpansiveness and the triangle inequality
also give, by induction up the tree,
\begin{equation}
 \norm{\Phi_k((H_v)) - \Phi_k((H'_v))}
 \le \sum_{v\in L_k}W_v\norm{H_v-H'_v}.
 \label{eq:tree-sensitivity}
\end{equation}
Indeed, at each parent the difference is at most the weighted sum of
the child differences.  Iterating multiplies the weights along each
path.  Consequently the root error is at most the largest cutoff
error, and changing only the output at $v\in L_k$ changes the root
by at most $W_v$ times the change at $v$.

\subsection{Predictions using earlier candidates}

We first describe the information associated with a candidate node
$v$ of depth $k\ge1$ during construction.  At this stage, all levels
of depth less than $k$, including their maps and weights, have
already been fixed.  Candidates at depth $k$ are selected one parent
group at a time, from left to right, and within each group in their
order of selection.  Let $I_v$ be the finite list of earlier
depth-$k$ candidates.  We retain their points $y_h$ and parent labels
even if some of these candidates have already been discarded from
the tree.  The map $G_v$ and point $y_v$ have not yet been chosen.

The later switching experiment replaces leaf outputs from right to
left. If $v$ is retained, then just after its block has switched, the \emph{ideal}
cutoff state uses $y_h$ at nodes strictly left of $v$ and a common
point $b$ at $v$ and to its right. The actual subtree outputs will
approximate these prescribed values. We must describe the root
outputs of this ideal state before the current level's weights and
later children are known.

Put
\[
 \Theta_v=\left\{(\theta_h)_{h\in I_v}:
    \theta_h\ge0,\quad
    \sum_{\substack{h\in I_v\\\widehat h=p}}\theta_h\le1
       \quad(p\in L_{k-1})\right\}.
\]
This is a compact subset of a finite-dimensional space.  If $I_v$
is empty, it consists of the single empty tuple.  For $b\in K$ and
$\theta\in\Theta_v$, set
\begin{equation}
 \begin{split}
 a_{v,p}(b;\theta)
   &=b+\sum_{\substack{h\in I_v\\\widehat h=p}}
                     \theta_h(y_h-b),\qquad p\in L_{k-1},\\
 Q_v(b;\theta)
   &=\Phi_{k-1}\bigl((G_p(a_{v,p}(b;\theta)))_{p\in L_{k-1}}\bigr).
 \end{split}
 \label{eq:prediction-definition}
\end{equation}
Each $a_{v,p}$ is a convex combination of $b$ and the earlier
candidates belonging to $p$, so it lies in $K$.  The residual
coefficient of $b$ represents all children whose outputs are to be
replaced by that common point.  Only maps and weights at depths less
than $k$ enter the evaluation of $Q_v$.

For fixed $v$, the map $\theta\mapsto Q_v(b;\theta)$ is norm
Lipschitz, uniformly over $b\in K$.  In fact, since $\diam K=1$,
\[
 \norm{a_{v,p}(b;\theta)-a_{v,p}(b;\eta)}
 \le \sum_{\substack{h\in I_v\\\widehat h=p}}
                      \abs{\theta_h-\eta_h}.
\]
Applying the nonexpansive maps $G_p$ and
Equation~\eqref{eq:tree-sensitivity} gives
\begin{equation}
 \norm{Q_v(b;\theta)-Q_v(b;\eta)}
 \le \sum_{h\in I_v}\abs{\theta_h-\eta_h}.
 \label{eq:prediction-lipschitz}
\end{equation}
For each $\theta$, weak compactness of $K$ gives the limit
\[
 Z_v(\theta)=\operatorname*{w-lim}_{j\to\mathcal U}
                      Q_v(b_j;\theta)\in K.
\]
The difference of the two sequences at $\theta$ and $\eta$
converges weakly along $\mathcal U$ to
$Z_v(\theta)-Z_v(\eta)$.  The norm ball of radius
$\sum_{h\in I_v}\abs{\theta_h-\eta_h}$ is weakly closed, so
Equation~\eqref{eq:prediction-lipschitz} implies the same bound for
$\norm{Z_v(\theta)-Z_v(\eta)}$.  Thus $Z_v$ is continuous into the
norm topology, and
\begin{equation}
 E_v=\{z_1,\ldots,z_k\}\cup Z_v(\Theta_v)
 \label{eq:tree-detector-set}
\end{equation}
is a nonempty norm-compact subset of $K$.  It is completely
determined before $G_v$ is selected. The dense test points will force
the branch functionals to tend to zero on the differences $z_i-z_j$,
and hence on $Y$; the compact prediction image provides comparison
limits for block detection. This argument takes weak limits
of the already evaluated root outputs; it does not pass a weak
limit through a nonlinear map.

\subsection{The inductive construction}

\begin{proposition}[Tree data]
\label{prop:tree-data}
There is an infinite ordered tree with the maps, points, and positive
weights described above.  Each nonroot node $v$ also has a budget
$e_v>0$, a functional $f_v\in X^*$ with $\norm{f_v}\le1$, and the
finite earlier-candidate list and predictions just defined.  These
data have the following properties:
\begin{enumerate}
 \item The budgets satisfy
 \begin{equation}
   \sum_{v\ne r}e_v\le\frac1{128}.
   \label{eq:tree-budget}
 \end{equation}
 Define
 \begin{equation}
  \tau_0=\frac12,\qquad
  \tau_k=2^{-k-1}
       \min\{e_u:u\in L_i,\ 1\le i\le k\}\quad(k\ge1).
  \label{eq:tree-tolerance}
 \end{equation}
 Every node $u\in L_k$ satisfies
 \begin{equation}
  \norm{\sum_{\widehat v=u}\alpha_v y_v-x}\le\tau_k.
  \label{eq:child-average}
 \end{equation}
 \item For every nonroot node $v$,
 \begin{equation}
  \sum_{p\prec v}\norm{G_p(y_v)-y_v}\le e_v,
  \qquad
  f_v(y_v-z)\ge1-e_v\quad(z\in E_v).
  \label{eq:tree-conditions}
 \end{equation}
 In particular, the second inequality holds at every $z_i$ with
 $i\le k$ when $v\in L_k$.
\end{enumerate}
The budgets, functionals, and prediction sets assigned to retained
nodes remain fixed when the remaining weights at their level and all
later levels are chosen.
\end{proposition}

\begin{proof}
Start with the root and $\tau_0=1/2$.  Number all candidate
selections throughout the construction by $\ell=1,2,\ldots$,
including candidates later discarded.  At selection number $\ell$,
assign the budget
\[
 e_v=2^{-\ell-8}
\]
before choosing the candidate's map.  The sum of all such budgets
is at most $1/256$, which is more than sufficient for
Equation~\eqref{eq:tree-budget} after discarding any candidates.

Suppose the finite levels through $L_{k-1}$ have been constructed
and $\tau_{k-1}$ has been fixed.  Visit the parents in $L_{k-1}$
from left to right.  For a given parent, apply
Lemma~\ref{lem:adaptive-selection} with tolerance $\tau_{k-1}$.
Its successive map choices will be the candidates for this parent's
children.  We now specify the threshold to be imposed at every
one of these choices.

Let $v$ be the next candidate.  Its budget has just been assigned.
Using all earlier candidates at the present level, form $I_v$,
$\Theta_v$, $Q_v$, and $E_v$ as above.  The ancestor maps are fixed,
so the number
\[
 C_v=\sum_{p\prec v}c_{G_p}
\]
is finite; take $c_{G_r}=0$.  By
Lemma~\ref{lem:compact-detector}, there is $\delta_v>0$ such that
any $y\in K$ with $\norm{Ty-y}<\delta_v$ admits a functional
$f\in X^*$, $\norm{f}\le1$, satisfying
$f(y-z)\ge1-e_v$ for every $z\in E_v$.  Choose the integer
threshold $m$ so large that
\[
  \frac1m<\delta_v,\qquad \frac{C_v}{m}\le e_v.
\]
This choice uses only data already available, and is made before
selecting $G_v$.

Require the next map in the adaptive selection to belong to
$\mathcal H_m$.  Every such choice has
$\norm{Ty_v-y_v}\le1/m$ by Equation~\eqref{eq:tail-defect}.
Equation~\eqref{eq:map-displacement} then gives
\[
 \sum_{p\prec v}\norm{G_p(y_v)-y_v}
 \le C_v\norm{Ty_v-y_v}\le e_v.
\]
After selecting $G_v$, choose the functional $f_v$ given by the
compact detector.  Both requirements in
Equation~\eqref{eq:tree-conditions} now hold.  In particular, they
hold for every map choice allowed by this step's threshold, so
they are compatible with the weak-neighborhood choices used in
Lemma~\ref{lem:adaptive-selection}.

That lemma terminates after finitely many candidates for the
parent under consideration.  Only then choose the convex
coefficients whose average of the points $y_v$ lies within
$\tau_{k-1}$ of $x$.  Discard the candidates with zero coefficient
and retain the others, in their original order, as this parent's
children; their positive coefficients are the weights $\alpha_v$.
The retained list is nonempty, and its weights sum to one.  Keep
the labels of discarded candidates in the finite lists used to
form subsequent predictions.  Thus no already defined prediction
or functional changes when zero coefficients are discarded.

There are only finitely many parents, so this completes a finite,
nonempty level $L_k$.  Every parent in $L_{k-1}$ now satisfies
Equation~\eqref{eq:child-average}.  Define $\tau_k$ by
Equation~\eqref{eq:tree-tolerance}; its defining minimum is over
a nonempty finite set of positive numbers.  Proceed to the next
level.  This constructs the entire tree and proves the stated
properties.
\end{proof}

Once level $L_k$ is complete, each retained $v\in L_k$ has a
particular parameter $\theta^v\in\Theta_v$: for $h\in I_v$, set
$\theta^v_h=\alpha_h$ if $h$ was retained, and set
$\theta^v_h=0$ otherwise.  The parent-group sums are at most one,
as required.  Write $h<_k v$ for the left-to-right order on
the retained level.  For $b\in K$, prescribe the ideal cutoff
outputs
\[
 \widetilde H_h^{v,b}=
 \begin{cases}
   y_h,&h<_k v,\\
   b,&h=v\ \text{or}\ v<_k h.
 \end{cases}
\]
For every parent $p\in L_{k-1}$, their weighted average is exactly
\[
 \sum_{\widehat h=p}\alpha_h\widetilde H_h^{v,b}
   =b+\sum_{\substack{h\in I_v\\\widehat h=p}}
                         \theta^v_h(y_h-b).
\]
Consequently
\begin{equation}
 \Phi_k((\widetilde H_h^{v,b})_{h\in L_k})
       =Q_v(b;\theta^v).
 \label{eq:prediction-cutoff}
\end{equation}
In a parent group earlier than that of $v$, all retained children
contribute their $y_h$ and their weights sum to one.  In the group
containing $v$, the residual coefficient of $b$ collects $v$ and
all later children.  In every later parent group, all ideal outputs
are $b$, so their average is $b$ regardless of how many children
were eventually selected or what their weights are.  This explains
why the prediction required no information about those future
children. Figure~\ref{fig:cutoff} illustrates this cutoff state.
The actual parameters $\theta^v$ are used only after
the weights are fixed; the compact sets $E_v$ were chosen to cover
all their possible values.


\begin{figure}[t]
\centering
\begin{tikzpicture}[>=Stealth,
  point/.style={circle,draw,minimum size=6mm,inner sep=1pt},
  untouched/.style={point,fill=white},
  switched/.style={point,fill=linkblue!15},
  every node/.style={font=\small}]
\node (root) at (0,2.4) {$r$};
\node (p1) at (-3,1.35) {$p_1$};
\node (p2) at (0,1.35) {$p_2$};
\node (p3) at (3,1.35) {$p_3$};
\draw (root)--(p1) (root)--(p2) (root)--(p3);
\node[untouched] (h1) at (-3.6,0) {$y_{h_1}$};
\node[untouched] (h2) at (-2.4,0) {$y_{h_2}$};
\node[untouched] (h3) at (-.6,0) {$y_{h_3}$};
\node[switched] (v) at (.6,0) {$b$};
\node[switched] (h5) at (2.4,0) {$b$};
\node[switched] (h6) at (3.6,0) {$b$};
\draw (p1)--(h1) (p1)--(h2) (p2)--(h3) (p2)--(v) (p3)--(h5) (p3)--(h6);
\node[below=2mm of v] {$v$};
\node[anchor=east] at (-4.3,0) {cutoff $L_k$};
\draw[->,thick] (3.6,-1.0)--(-3.6,-1.0)
 node[midway,below] {order of switching};
\end{tikzpicture}

\caption{An ideal cutoff state immediately after the block below $v$
has switched. Earlier nodes retain their labels; $v$ and the nodes to
its right have output $b$. Each parent applies its map to the weighted
average of its children. The diagram displays a two-level example of
this rule; in a deeper tree the same evaluation continues up to the root.
The coefficients on the earlier labels specify $\theta^v$, and their
residual mass multiplies $b$.}
\label{fig:cutoff}
\end{figure}
